\section{The dual approximation}
\label{sec:dual}

\subsection{The relaxation and why it bounds}

\citet[Section~2.3]{KWG2011} obtain a lower bound on a minimization by dualizing the equalities $Ax + s = b$ of the problem with slacks,
restricting the dual multipliers to linear rules $y(\xi) = Y\xi$, and carrying out the maximization over $Y$.
What remains is the original problem with the equalities required to hold only in expectation against $\xi$,
$\E[(Ax(\xi) + s(\xi) - b(\xi))\xi^\T] = 0$, and $s(\xi) \ge 0$ almost surely.
Its tractable form \citep[(2.8) and (4.6)]{KWG2011} replaces the policies by their projections $X$, $S$ on the linear functions of $\xi$,
defined by $X\Mom = \E[x(\xi)\xi^\T]$ and $S\Mom = \E[s(\xi)\xi^\T]$,
and the almost sure nonnegativity of $s$ by a finite set of linear inequalities on $S\Mom$, one per row of $W$.

We write the test functions as a cone:
let $\Wp$ be the matrix with the rows $\eone^\T$ and $(-h_{bp}, W_{bp})$ for every row $p$ of every support block $b$, embedded in the $k$ columns of $\xi$,
so that $\Wp\xi = (\xi_1,\ W\eta - h\xi_1)$ and
\[
\xi \in \Xis \implies \Wp\xi \ge 0.
\]
In \citet[(2.1b)]{KWG2011}, $W = (\eone, -\eone, \hat W^\T)^\T$ and $h = (1, -1, 0, \dots, 0)$, so that $W - h\eone^\T$ consists of two zero rows and the rows $\hat W$;
the rows of $\Wp$ other than $\eone^\T$ play the role of $\hat W$, with the constants of $h$ moved to the column of $\xi_1$.
The dual approximation solved by the package is
\begin{equation}
\label{eq:DLDR}
\tag{D-LDR}
\begin{aligned}
\opt_{X \in \Rules} \quad & \text{objective of \cref{prop:objective}} \\
\text{s.t.} \quad & A_e X = B_e, \\
& S_r = \sigma_r(b_r^\T - a_r^\T X) \\
& S_{r1} \ge 0 \text{ for the deterministic } r \in \Rows,\\
& S_r \Mom \Wp^\T \ge 0 \quad \text{for the other } r \in \Rows.
\end{aligned}
\end{equation}
The inequalities $S_r\Mom\Wp^\T \ge 0$ are $\E[s_r(\xi)\,(\Wp\xi)^\T] \ge 0$ for linear $s_r$:
the slack must be nonnegative ``on average against every function that is nonnegative on $\Xis$ by construction''.

\begin{theorem}[Dual bound]
\label{thm:dual-bound}
Suppose \cref{ass:support} and the following conditions hold.
\begin{enumerate}[label=(\roman*)]
\item\label{it:lce} \emph{Linear conditional expectations}: for every stage $t$ there is a matrix $N_t$ such that
$\E[\xi \mid \xi_{1:\nvis(t)}] = N_t\, \xi_{1:\nvis(t)}$ almost surely.
\item\label{it:convex} $P \succeq 0$ if $\opt = \min$, and $P \preceq 0$ if $\opt = \max$.
\end{enumerate}
Then for every feasible policy $x$ of \eqref{eq:SP} with square-integrable components, the matrix $X = \E[x(\xi)\xi^\T]\Mom^{-1}$ is feasible in \eqref{eq:DLDR},
and its objective is at least as good as that of $x$.
Consequently, for a minimization, $\val\eqref{eq:DLDR} \le \val\eqref{eq:SP} \le \val\eqref{eq:PLDR}$, and the reverse for a maximization.
\end{theorem}

\begin{proof}[Proof sketch]
The equalities hold almost surely, so in expectation against $\xi^\T$, which gives $A_eX\Mom = B_e\Mom$, hence $A_eX = B_e$.
For an inequality, $s_r(\xi) \ge 0$ and $\Wp\xi \ge 0$ almost surely, so $\E[s_r(\xi)(\Wp\xi)^\T] \ge 0$,
and $\E[s_r(\xi)\xi^\T] = S_r\Mom$ with $S_r = \sigma_r(b_r^\T - a_r^\T X)$ by linearity.
The objective is linear in $\E[x\xi^\T] = X\Mom$, except for the term in $P$:
writing $x = X\xi + \varepsilon$ with $\E[\varepsilon\xi^\T] = 0$, the cross term vanishes and $\E[\varepsilon^\T P \varepsilon]$ has the sign of $P$.
It remains to show that $X \in \Rules$, which is where condition~\ref{it:lce} enters:
the projection of a function of the history $\xi_{1:\nvis_i}$ on the linear functions of the whole of $\xi$ only involves the history if the conditional expectation of the future given the history is linear,
which is \citet[Lemma~2]{KWG2011}.
An integer decision is a constant, and its projection is the same constant, so integrality is preserved.
The complete argument is in \cref{app:dual}.
\end{proof}

This is the lower bound of \citet[Section~4]{KWG2011}, whose tractable form is their Theorem~3, under their condition $\E_t[\xi] = M_tP_t\xi$,
with a quadratic objective and integer first-stage decisions added.
They also show that \eqref{eq:DLDR} is exact for their relaxed problem when the original one is strictly feasible \citep[Proposition~4]{KWG2011};
the bound itself does not need strict feasibility.
\citet{GWK2015} assume instead that the stages are independent \citep[condition (M4)]{GWK2015},
which implies condition~\ref{it:lce} \citep[footnote~3]{KWG2011}.

% What condition (i) implies?  In this "no stage for constraints", is it needed?
\begin{remark}[What condition~\ref{it:lce} asks of the laws]
\label{rem:lce}
Distinct declarations are independent, so condition~\ref{it:lce} holds as soon as each declaration lies in one stage,
and more generally when the law of each declaration has linear conditional expectations across its own stages.
A multivariate normal law has this property, and so has its truncation to a confidence ellipsoid (\cref{sec:confidence-normal});
a discrete joint law of several stages generally does not.
The package checks condition~\ref{it:lce} law by law, for the laws whose components are in several stages (\code{\_linear\_conditional\_expectations}, \file{src/probability/moments.jl}):
it holds for a product of univariate laws (independent components), a normal law and \code{ConfidenceMvNormal};
for an \code{MvDiscreteNonParametric} it is checked exactly, for each stage $t$, by fitting an affine function of the history to the conditional means of the later components, one per distinct history, which must leave no residual;
for any other law, such as a law conditioned by rejection, it is unknown.
When it fails or is unknown for some law, the package warns, and marks the bound as not certified (\code{\_stage\_dependent\_laws}, \file{src/formulations/dual.jl}).
Correlation is not the test: a correlated law may satisfy condition~\ref{it:lce}, such as the truncated normal law above,
and an uncorrelated one may not, such as $\eta_2 = \eta_1^2$ with $\eta_1$ uniform on $\{-1, 0, 1\}$.
\end{remark}

\begin{remark}[The row $\eone^\T$]
\label{rem:e1-row}
In \citet{KWG2011}, $\Xis$ is compact, and the cone $\{z : (W - h\eone^\T)z \ge 0\}$ already implies $z_1 \ge 0$ \citep[proof of Proposition~3]{KWG2011};
the rows $\pm\eone^\T$ of their $W$ become zero rows of $W - h\eone^\T$.
The package allows unbounded blocks, and keeps the row $\eone^\T$ of $\Wp$ explicitly,
that is, the inequality $\E[s_r(\xi)] \ge 0$.
Without it, a model in which no uncertainty is bounded would have no inequality left in its dual approximation;
the previous version of the package had this defect.
\end{remark}

The relaxation depends on the set $\Xis$, not on the inequalities that describe it.

\begin{proposition}[Invariance under the description of $\Xis$]
\label{prop:description}
Adding to $W\eta \ge h$ an inequality that is valid on $\Xis$ does not change the feasible set of \eqref{eq:DLDR}.
In particular, redundant rows and positive scalings of rows have no effect on the dual bound.
\end{proposition}

\begin{proof}
Let $w^\T\eta \ge g$ be valid on the nonempty polyhedron $\{W\eta \ge h\}$.
By the affine form of Farkas' lemma there are $c \ge 0$ and $c_0 \ge 0$ with $w = W^\T c$ and $g = h^\T c - c_0$.
Then $w^\T\eta - g\xi_1 = c^\T(W\eta - h\xi_1) + c_0\xi_1$ is a nonnegative combination of the rows of $\Wp$,
and the new inequality $S_r\Mom(-g, w)^\T \ge 0$ is the same combination of inequalities of \eqref{eq:DLDR}.
\end{proof}

In the language of \citet[Proposition~3]{KWG2011}, the cone defined by the rows of $\Wp$ is the closed convex cone generated by $\Xis$.
The row $\eone^\T$ is needed for the proposition when $\Xis$ is unbounded.

\subsection{The structured form}

Written literally, $S_r\Mom\Wp^\T$ is a row with one entry per row of $\Wp$, each a sum of $k$ products of an entry of $S_r$, itself an expression of $X$, by an entry of the dense matrix $\Mom\Wp^\T$.
This is what made the dual of the previous version explode.
Substituting \eqref{eq:structured-moments},
\begin{equation}
\label{eq:dual-identity}
S_r \Mom \Wp^\T = (S_r\mub)(\Wp\mub)^\T + S_r\Cv\Wp^\T.
\end{equation}
Introduce the scalar $u_r \defeq S_r\mub$, the expected slack.
The row $\eone^\T$ of $\Wp$ gives $u_r \ge 0$, since $\Cv\eone = 0$ and $\eone^\T\mub = 1$.
The row $p$ of support block $b$ gives
\begin{equation}
\label{eq:dual-cone-row}
u_r\,(W_{bp}\mu_b - h_{bp}) + S_r\, \Cv_{:,b} W_{bp}^\T \ge 0,
\end{equation}
where $W_{bp}\mu_b - h_{bp} \ge 0$ because the mean lies in $\Xib{b}$,
and $\Cv_{:,b}$ are the columns of $\Cv$ of the coordinates of block $b$.
The matrix $V_b \defeq \Cv_{:,b}W_b^\T$ is computed once per support block (\code{\_dual\_cones}),
and its row $j$ is nonzero only if column $j$ shares a covariance block with some coordinate of $b$.
In \eqref{eq:dual-cone-row}, $S_r$ has nonzeros only in $\Jr_r$, so:
\begin{itemize}[nosep]
\item only the support blocks that share a covariance block with a structural column of $r$ give inequalities with a slack term;
\item a row of $\Wp$ whose term $S_r V_{b,:p}$ is structurally zero reduces to $u_r(W_{bp}\mu_b - h_{bp}) \ge 0$, implied by $u_r \ge 0$, and is dropped:
all the blocks that the inequality does not see \emph{collapse} into the single inequality $u_r \ge 0$;
\item each remaining inequality has one term in $u_r$ and one term per structural column $j$ with $(V_b)_{jp} \neq 0$.
\end{itemize}
The package creates, for each inequality $r$ that is not deterministic, the variables $u_r \ge 0$ and $S_{rj}$, $j \in \Jr_r$,
with the equalities $S_{rj} = \sigma_r(b_{rj} - (a_r^\T X)_j)$ and $u_r = \sum_{j\in\Jr_r} \mub_j S_{rj}$,
and the inequalities \eqref{eq:dual-cone-row} that have a slack term (\code{\_add\_dual\_cone!}).
Restricting $S_r$ to $\Jr_r$ is not a restriction: the other entries are zero for every $X \in \Rules$, as noted after \eqref{eq:slack-rule}.
This is the masking of $S$ by the history that \citet[(4.6)]{KWG2011} write as $S_tP_t$.

\begin{proposition}[Structured dual]
\label{prop:structured-dual}
The structured problem above and the dense problem \eqref{eq:DLDR} have the same feasible rules $X$ and the same objective.
\end{proposition}

\begin{proof}
By \eqref{eq:dual-identity}, the inequalities $S_r\Mom\Wp^\T \ge 0$ are $u_r \ge 0$ and the inequalities \eqref{eq:dual-cone-row} for all rows of all blocks, with $u_r = S_r\mub$.
Those that are dropped are implied by $u_r \ge 0$ and $W_{bp}\mu_b - h_{bp} \ge 0$.
The variables $u_r$ and $S_{rj}$ are determined by $X$ through the equalities.
\end{proof}

The previous version's spike validated this identity numerically on the hydrothermal problem
(objectives equal to a relative difference of $10^{-12}$ on 12 instances),
and the test suite compares both formulations with dense oracles (\file{test/oracles/dense.jl}).
The auxiliary $u_r$ keeps each inequality \eqref{eq:dual-cone-row} as sparse as the covariance allows:
substituting it out would add the $\abs{\Jr_r}$ terms $(\Wp\mub)_p\mub_j S_{rj}$ to every inequality of the row.
Whether to keep it is a question of solver behaviour more than of size, discussed in \cref{sec:computation}.

\subsection{Reporting the bound}

The rule $X$ of \eqref{eq:DLDR} is not a policy: it need not be feasible for any outcome.
The package reports only its value, as the objective bound of the model, and the relative gap between it and the primal value.
With integer decisions, the inner problem is a mixed-integer program, and the bound is the inner solver's bound, which needs no incumbent;
for a linear program, it is the optimal value, or the dual objective value of the inner solver if the solve stopped early.
Since \eqref{eq:DLDR} is a relaxation, its infeasibility proves the infeasibility of \eqref{eq:SP} (under the conditions of \cref{thm:dual-bound}),
while an infeasible primal approximation only shows that no linear rule is feasible:
the package reports \code{INFEASIBLE} only on the evidence of the dual, and \code{LOCALLY\_INFEASIBLE} otherwise.
The bound is \emph{certified} when the moments are exact, $\Mom \succ 0$, and every law spanning several stages is known to satisfy condition~\ref{it:lce} (\cref{rem:lce});
otherwise the result says why it is not (\cref{sec:rejection}).
